\section{Discussion}
\label{sec:discussion}

The result we did not expect is that the sequence-level optimum anti-concentrates the
quantity it constrains. Stated after the fact it is almost obvious, since an expectation
constraint says nothing about individual draws and the cheapest way to move a mean is to move
the tails. It nonetheless inverts the usual reading of a per-step surrogate as a compromise. The distinction between constraining an expectation and constraining every
trajectory is old; what we add is a measurement of how far apart those
two objects are for a concrete generative model, and the observation that the per-step route
lands closer to the second.

The boundaries are as informative as the results. The steering ceiling says the
projection reweights rather than creates. The composition asymmetry says a running budget
is a remaining quantity, so constraints that want more of it fail differently from
constraints that want less. The determinability result says the expressible constraint
class is chosen when the serialization is chosen, which is a design decision usually made
for other reasons. And best-of-$n$'s accuracy win says the contribution is cost structure,
not accuracy: a claim weaker than the one we set out to make and better supported than it
would have been. \Cref{si:defects} records three errors caught during
this work and the mechanism that caught each, since each would otherwise have produced a
plausible-looking wrong number.

%%%%%%%%%%%%%%%%%%%%%%%%%%%%%%%%%%%%%%%%%%%%%%%%%%%%%%%%%%%%%%%%%%%%%%%%%%%%%%%%
